\chapter{Dentistry}

\medterm{Amalgam}-- A dental restorative material composed of elemental mercury (50%) combined with a metallic alloy powder containing silver, tin, copper, and zinc (most modern: high-copper amalgam). Advantages: durable (lasts 10-15+ years), low cost, moisture-tolerant (can be placed in areas not perfectly dry), longest track record of any restorative material. Disadvantages: esthetic (silver color, not tooth-colored), requires mechanical retention (undercuts, slots), tooth structure removal for retention, and concern over mercury content (FDA: safe for most patients; avoids placement/removal in pregnant women and those with mercury allergy). Placement: cavity preparation with undercuts, mixing amalgam in a capsule (trituration), condensation into the cavity, and carving to contour. Mercury exposure: dental personnel may have higher urine mercury but below toxic thresholds. Amalgam separators (capture >95% of amalgam waste) are required by the EPA. Alternatives: composite resin (tooth-colored), glass ionomer, and ceramic. Use has declined in many countries (Sweden, Norway phased out for environmental reasons).

\medterm{Apical Periodontitis}-- Inflammation and destruction of the periradicular tissues (periodontal ligament and alveolar bone at the apex) caused by bacterial infection of the root canal system (necrotic pulp). Presents with pain on biting/percussion (symptomatic) or asymptomatic with radiographic periapical radiolucency (asymptomatic). Acute apical abscess: severe pain, swelling (facial cellulitis, space infection), fever, and malaise. Chronic apical periodontitis: asymptomatic, well-defined periapical radiolucency. Phoenix abscess: acute exacerbation of chronic lesion. Treatment: root canal therapy (remove infected pulp, clean, shape, disinfect, and obturate the root canal system). Incision and drainage for fluctuant swelling (through the gingiva or mucosa). Antibiotics: for systemic signs, facial swelling, or immunocompromised patients (amoxicillin, clindamycin if penicillin-allergic).

\medterm{Apicoectomy (Root-End Resection)}-- A surgical endodontic procedure performed when conventional root canal therapy or retreatment is unsuccessful due to persistent apical disease. Indications: persistent apical periodontitis after RCT, obstructed canal (calcification, separated instrument, broken file), apical perforation, root-end cyst, and biopsy of periapical lesion. Procedure: local anesthesia, full-thickness mucoperiosteal flap elevation, osteotomy (bone window) to access the root apex, apical root-end resection (3 mm of apex removed), ultrasonic root-end preparation (3 mm deep), root-end filling (MTA mineral trioxide aggregate or Super-EBA), and flap repositioning/suturing. Success rate: 75-90% with MTA apical seal. Post-operative: mild to moderate pain, swelling, avoid heavy chewing for 1 week.

\medterm{Bridge (Fixed Partial Denture)}-- A fixed dental prosthesis that replaces one or more missing teeth by attaching to adjacent natural teeth (abutments) or implants. Typical: a three-unit bridge (two abutments, one pontic — artificial tooth replacing the missing tooth). Types: traditional (full-coverage crowns on abutments, pontic in between), cantilever (pontic supported on only one side — more conservative but higher failure rate), resin-bonded/Maryland bridge (wings bonded to the lingual surfaces of adjacent teeth — minimal tooth reduction, for single anterior tooth replacement), and implant-supported bridge (most conservative of adjacent teeth — no preparation of healthy adjacent teeth). Pontic design: modified ridge-lap or ovate pontic (hygienic, esthetic). Materials: PFM, full ceramic (lithium disilicate, zirconia), and metal. Relative contraindications: long span, heavy occlusal forces (bruxism), short clinical crowns (insufficient retention). Longevity: 10-15 years. Failure: abutment tooth fracture/decay (most common, 10-20% at 10 years), debonding, and pontic fracture.

\medterm{Bruxism}-- A repetitive jaw-muscle activity characterized by clenching or grinding of the teeth, occurring during sleep (sleep bruxism) or wakefulness (awake bruxism). Prevalence: 8-15% of adults, up to 30% of children. Etiology: multifactorial — psychosocial stress, anxiety, sleep disorders (OSA, snoring), occlusal interference, genetics, caffeine, alcohol, smoking, and certain medications (SSRIs have been implicated — may cause or worsen). Clinical consequences: tooth wear/attrition (incisal and occlusal surfaces — flattened cusps, exposed dentin), tooth sensitivity, fractures of restorations/cusps, masticatory muscle hypertrophy (masseter hypertrophy visible), TMJ pain, morning headache, and tongue/cheek scalloping. Diagnosis: clinical exam (wear facets), history (reports of grinding from bed partner), and polysomnography (for sleep bruxism, definitive, shows rhythmic masticatory muscle activity RMMA). Treatment: occlusal splint/night guard (hard acrylic or soft, reduces wear and muscle activity), stress management, biofeedback, sleep hygiene, and in severe cases, botulinum toxin injection into masseter and temporalis muscles. Pharmacotherapy: clonidine, clonazepam (low-dose, short-term). No cure; management is symptomatic.

\medterm{Composite Resin (Tooth-Colored Filling)}-- An esthetic dental restorative material composed of a resin matrix (Bis-GMA, UDMA, TEGDMA) filled with inorganic particles (silica, barium glass, zirconia). Categories by filler size: microfilled (0.04 µm, polishable but weak, for anterior), hybrid (0.5-1 µm + microfillers, universal), nanofilled/nanohybrid (5-100 nm nanoparticles, best properties). Advantages: esthetic (color matched to the tooth), minimal tooth preparation (bonded restoration, preservation of sound tooth structure), and repairable. Application: acid-etch (37% phosphoric acid for 15-30 seconds — removes smear layer, creates microporosity for micromechanical bond), apply bonding agent (primer + adhesive), place composite in increments (<2 mm each to reduce polymerization shrinkage), light-cure (LED light 400-500 nm for 10-40 seconds), finish and polish. Polymerization shrinkage: 1.5-3%, can cause marginal gaps, post-operative sensitivity, microleakage, and secondary caries. Bulk-fill composites (can be placed in 4-5 mm increments) reduce technique sensitivity. C-factor: ratio of bonded to unbonded surfaces; higher C-factor = more shrinkage stress. Longevity: 5-8 years in posterior teeth, longer in anterior.

\medterm{Crown (Dental Crown)}-- A full-coverage dental restoration that encases the entire visible portion of a tooth above the gum line, restoring its shape, size, strength, and appearance. Indications: large caries/cracks (insufficient tooth structure for a filling), crown fracture, after root canal therapy (posterior teeth — endodontically treated teeth are brittle), cosmetic enhancement, and as a bridge abutment. Crown materials: full metal (gold alloy — most durable, wear-resistant, minimal tooth reduction — 1-1.5 mm), PFM (porcelain fused to metal — good esthetics and strength, metal margin may show with gingival recession), all-ceramic (lithium disilicate — e.max, most esthetic, high strength, 1-1.5 mm reduction), and zirconia (yttria-stabilized zirconia — highest strength, used for posterior crowns and bridges, CAD/CAM fabricated, opaque). Preparation: 1-2 mm tooth reduction circumferentially with a shoulder or chamfer margin. Impression: digital (intraoral scanner — iTero, Trios, PrimeScan) or conventional (polyvinyl siloxane or alginate). Temporary crown: acrylic or bis-acryl composite temporary for 1-2 weeks while the permanent crown is fabricated. Cementation: resin cement (for ceramic and zirconia), glass ionomer, or zinc phosphate. Longevity: 10-15+ years, with proper hygiene and regular recall.

\medterm{Deep Sedation and General Anesthesia in Dentistry}-- Pharmacologically induced states of consciousness suppression for dental treatment, used for patients with severe dental phobia, special needs (developmental delay, autism, cerebral palsy), young children requiring extensive treatment, and lengthy/complex surgical procedures (impacted third molars, dental implants, orthognathic surgery). Practitioner: DDS/DMD with a dental anesthesia residency permit or certified registered nurse anesthetist (CRNA). Levels: moderate conscious sedation (IV: midazolam, fentanyl; oral: triazolam, hydroxyzine — patient can maintain airway and respond to verbal commands), deep sedation (propofol, ketamine — purposeful response to painful stimulation, may require airway support), and general anesthesia (endotracheal intubation, controlled ventilation). Monitoring: pulse oximetry, capnography, ECG, NIBP. Patient selection: ASA I-II generally safe, ASA III-IV require medical consultation. Risks: respiratory depression (opioids, propofol), laryngospasm, aspiration, hypotension, cardiac arrhythmia. Recovery: monitored PACU with Aldrete scoring.

\medterm{Dental Abscess}-- A localized collection of pus in the oral cavity, classified by origin: periapical (from necrotic pulp — most common), periodontal (from periodontal pocket), and pericoronal (pericoronitis around partially erupted tooth, most commonly lower third molar). Signs: severe, throbbing pain, tooth tender to percussion (periapical), swelling (intraoral, extraoral, facial), cellulitis (diffuse swelling, systemic involvement — Ludwig angina for submandibular space infection, life-threatening airway compromise), lymphadenopathy, fever, and malaise. Radiographic: periapical radiolucency (periapical abscess may not be visible early — bone demineralization >30% needed for radiographic visibility). Treatment: establish drainage (root canal therapy — access the pulp chamber and allow drainage through the canal, or incision and drainage for fluctuant swelling), removal of the cause (RCT, extraction), and antibiotics for systemic involvement (amoxicillin 500 mg TID x 5-7 days, clindamycin 300 mg QID if penicillin-allergic). Severe infections: IV antibiotics, hospital admission, and surgical drainage.

\medterm{Dental Calculus}-- See Tartar (Dental Calculus).

\medterm{Dental Caries (Tooth Decay)}-- A chronic, progressive, multifactorial disease characterized by demineralization of tooth enamel and dentin by organic acids produced by bacterial fermentation of dietary carbohydrates. Etiology: cariogenic bacteria (Streptococcus mutans, Lactobacillus, Actinomyces), fermentable carbohydrates (sucrose most cariogenic), susceptible tooth surface, and time. Caries classification: pit and fissure (occlusal surfaces), smooth surface (interproximal, buccal/lingual), root caries (exposed root surfaces in recession), and recurrent caries (around existing restorations). Depth: incipient/reversible (enamel only, white spot lesion), moderate (dentin involvement, may be reversible if <1/3 into dentin), deep (approaching pulp, risk of pulpitis), and pulpal involvement (pulpitis, necrosis). Diagnosis: visual-tactile exam (ICDAS criteria), bitewing radiographs (detect interproximal caries not visible clinically), transillumination (DiFOTI), and laser fluorescence (DIAGNOdent). Prevention: fluoride (toothpaste, varnish, water fluoridation), dietary counseling (limit sugar frequency), pit and fissure sealants, and regular recall visits. Treatment: minimally invasive (resin infiltration for incipient lesions, GI/temporary), direct restoration (composite resin, amalgam, glass ionomer), indirect restoration (inlay/onlay/crown), and root canal therapy for pulpal involvement. ICDAS (International Caries Detection and Assessment System): codes 0-6 based on visual appearance from sound to extensive cavitation.

\medterm{Dental Emergencies}--  severe toothache, knocked-out teeth, cracked or fractured teeth, abscesses, and uncontrolled bleeding from the mouth. A knocked-out permanent tooth should be placed back in the socket or in milk and seen by a dentist within 30 minutes for the best chance of reimplantation. Dental abscesses require prompt treatment to prevent spread of infection to adjacent tissues or the bloodstream.
\medterm{Dental Health}-- Dental health involves the care of teeth, gums, and supporting structures to prevent common problems like cavities, gum disease, and tooth loss. Regular brushing with fluoride toothpaste, flossing, and professional cleanings control plaque buildup and detect issues early. Poor dental health is linked to systemic conditions including cardiovascular disease, diabetes complications, and adverse pregnancy outcomes.

\medterm{Dental Public Health}-- The non-clinical dental specialty focused on preventing oral disease and promoting oral health at the population level through community-based programs, policy development, and research. Activities: community water fluoridation (optimal 0.7 mg/L, reduces caries by 25-40%), school-based sealant programs, fluoride varnish programs, oral health education, epidemiology of oral disease (NHANES data), and access-to-care initiatives. Key metrics: DMFT index (Decayed, Missing, Filled permanent Teeth), DMFS (surfaces), dmft for primary dentition. Public health challenges: disparities in oral health access (low-income, rural, racial/ethnic minorities), dental caries prevalence (most common chronic disease of childhood — 5x more common than asthma), and oral cancer screening.

\medterm{Dentin}-- The calcified tissue of the tooth that lies beneath the enamel (crown) and cementum (root), forming the bulk of the tooth. Dentin is composed of approximately 70% hydroxyapatite (mineral), 20% organic matrix (primarily type I collagen), and 10% water. It contains microscopic tubules (dentinal tubules) that radiate from the pulp chamber to the enamel--dentin junction, housing odontoblast processes and allowing transmission of sensory stimuli (heat, cold, pressure) to the pulp as pain. Dentin is produced by odontoblasts throughout life (secondary and tertiary dentin). When exposed by caries or attrition, dentin hypersensitivity may occur. Dentinogenesis imperfecta is a genetic disorder causing opalescent, weakened dentin.

\medterm{Dentition}-- The arrangement, number, and type of teeth in the dental arch at a given stage of life. Primary dentition (deciduous teeth, baby teeth): 20 teeth, erupting from ~6 months to ~2.5 years. Mixed dentition: phase when both primary and permanent teeth are present (~6--12 years). Permanent dentition (adult): 32 teeth including 8 incisors, 4 canines, 8 premolars, and 12 molars (including 4 third molars/wisdom teeth). Eruption timing varies but follows a predictable sequence. Dental formula records the number and type of teeth in each quadrant. Abnormalities: anodontia (congenital absence), supernumerary teeth (extra teeth), malocclusion, and impacted teeth.

\medterm{Denture}-- A removable dental prosthesis that replaces missing teeth and adjacent tissues. Complete denture: replaces all teeth in an arch (maxillary or mandibular). Immediate denture: placed immediately after extraction of remaining teeth (allows patient to never be without teeth; requires relining after bone resorption). Partial denture (RPD — removable partial denture): replaces some but not all teeth in an arch; held in place by clasps (cast metal — cobalt-chromium, or flexible nylon Valplast) and rests on abutment teeth. Implant-retained overdenture: denture that snaps onto 2-4 dental implants (greatly improved retention and stability, especially for the mandibular complete denture — 95% of lower denture problems solved by 2 implants). Implant-supported fixed complete denture: fixed prosthesis on 4-6 implants (All-on-4 concept). Adaptation period: 2-4 weeks for speech, eating, and saliva control. Denture stomatitis (denture sore mouth): Candida albicans infection under the denture — improved by removing the denture at night, proper cleaning, and antifungal (nystatin, clotrimazole). Complications: poor fit due to bone resorption (requires reline — hard or soft tissue conditioner), fracture (may require repair), and sore spots (adjustment with acrylic bur).

\medterm{Dry Socket (Alveolar Osteitis)}-- A common post-extraction complication characterized by severe, throbbing pain starting 2-3 days after tooth extraction, caused by premature loss of the blood clot from the extraction socket, exposing underlying bone. Risk factors: tobacco use (nicotine reduces blood supply, clot dissolution from sucking), traumatic extraction, oral contraceptives (estrogen increases fibrinolysis), poor oral hygiene, and mandibular teeth (especially third molars — higher density bone, less blood supply). Treatment: gentle socket irrigation with warm saline to clear debris, placement of a medicated dressing (Alvogyl — butamben, eugenol, iodoform; or eugenol-soaked ribbon gauze — zinc oxide-eugenol) for immediate pain relief, and NSAIDs. The dressing may need to be changed every 1-3 days for 5-10 days until granulation tissue covers the bone. Prevention: post-operative instructions (no smoking, no spitting, no drinking through straws for 48-72 hours).

\medterm{Endodontics}-- The dental specialty concerned with the morphology, physiology, and pathology of the dental pulp and periradicular tissues. Endodontists perform root canal therapy (pulpectomy, pulpotomy), root canal retreatment (for failed prior RCT), endodontic surgery (apicoectomy — root-end resection with retrograde filling), and management of traumatic dental injuries (crown fractures, root fractures, luxation, avulsion). Common diagnoses: irreversible pulpitis (severe, lingering pain to thermal stimuli), pulp necrosis (no response to cold, periapical radiolucency), and acute apical abscess (pain, swelling, drainage). Treatment involves cleaning and shaping the root canal system (hand files and rotary NiTi files), disinfection (sodium hypochlorite irrigation 1-6%, chlorhexidine, calcium hydroxide), and obturation (gutta-percha with sealer). Cone-beam CT (CBCT) is used for complex anatomy (extra canals, root resorption, fractures). Success rate: >90% for initial root canal treatment.

\medterm{General Dentistry}-- The primary dental care provider responsible for the diagnosis, prevention, and treatment of common oral diseases including dental caries (cavities), periodontal disease, and oral pathology. General dentists perform restorative procedures (fillings, crowns, bridges), simple extractions, root canal therapy, and preventive care (cleanings, fluoride, sealants). They also coordinate comprehensive treatment plans and refer to specialists when needed. Patients are typically seen every 6 months for routine recall examinations and prophylaxis.

\medterm{Gingival Hyperplasia (Overgrowth)}-- Abnormal enlargement of the gingiva. Drug-induced: phenytoin (antiepileptic — most severe, 50% of patients), cyclosporine (immunosuppressant, 30%), and calcium channel blockers (nifedipine, amlodipine — 20%). Other causes: hereditary gingival fibromatosis (rare, autosomal dominant), hormonal (pregnancy, puberty), inflammatory (chronic plaque-induced gingivitis), and leukemic infiltration (acute monocytic/myelomonocytic leukemia). Presents: firm, pink, lobulated enlargement of the interdental papillae and marginal gingiva, may cover the clinical crowns. Bleeding and secondary inflammation common when plaque control is poor. Treatment: improved oral hygiene (reduces inflammatory component, drug-induced overgrowth persists), chlorhexidine mouthwash, surgical excision (gingivectomy with scalpel, electrosurgery, or laser — CO2, diode). Discontinuation or substitution of the causative drug if possible (consult with prescribing physician). Phenytoin substitute: levetiracetam, lamotrigine. Recurrence is common if the drug continues.

\medterm{Gingivitis}-- Reversible inflammation of the gingiva (gums) caused by bacterial plaque biofilm accumulation at the gingival margin. Clinical signs: erythema, edema (blunted papilla), bleeding on probing (BOP — first sign, before visible redness), and loss of gingival stippling. No bone loss or clinical attachment loss (true gingivitis). Pregnancy gingivitis: hormone-mediated exaggerated response to plaque in the second/third trimester. Gingival hyperplasia/enlargement: phenytoin, cyclosporine, calcium channel blockers, and hormonal changes. Treatment: professional prophylaxis (scaling and polishing), improved oral hygiene (brushing twice daily, flossing, antigingivitis toothpaste — stannous fluoride, triclosan), and chlorhexidine mouthwash for short-term use. Resolution: 10-14 days with plaque control.

\medterm{Impacted Tooth}-- A tooth that fails to fully erupt into the normal functional position within the expected time frame, most commonly the third molars (wisdom teeth — 90% of impactions), followed by maxillary canines (2-3%), mandibular premolars, and maxillary central incisors (mesiodens, supernumerary). Third molar impaction classification: relationship to the ramus and second molar (Classes I, II, III), depth (Level A, B, C), and angulation (mesioangular most common, vertical, horizontal, distoangular). Indications for extraction: pericoronitis (recurrent infection of the operculum over partially erupted tooth), caries on the distal of the second molar, non-restorable caries, periodontal disease, root resorption, orthodontic reasons, cyst formation (dentigerous cyst), and prophylactic removal (controversial, NICE guidelines recommend against routine). Complications of extraction: dry socket (alveolar osteitis — most common, 1-5%, severe pain 2-3 days post-op, treat with irrigation and medicated dressing — Alvogyl, eugenol paste), inferior alveolar nerve injury (temporary 5-10%, permanent <1%), lingual nerve injury, maxillary sinus communication, mandible fracture (rare), and infection.

\medterm{Ludwig Angina}-- A severe, rapidly spreading cellulitis of the sublingual, submandibular, and submental spaces, most commonly from an odontogenic source (lower second/third molar infection). Life-threatening due to airway compromise from massive bilateral swelling elevating and posteriorly displacing the tongue (airway obstruction). Signs: bilateral submandibular swelling (bull neck appearance), elevation of the tongue, trismus, dysphagia, drooling, fever, and impending airway obstruction. Pathogens: Streptococcus viridans, anaerobes (Peptostreptococcus, Bacteroides, Fusobacterium), and mixed oral flora. Airway management is the priority: awake fiberoptic intubation or cricothyrotomy/tracheostomy (may be extremely difficult — early intervention before deterioration). Empiric IV antibiotics: high-dose penicillin + metronidazole (or ampicillin-sulbactam, clindamycin). Surgical drainage: bilateral submandibular and submental incisions with placement of Penrose drains. Extraction of the causative tooth. ICU monitoring.

\medterm{Malocclusion}-- Misalignment or incorrect positioning of the teeth when the jaws are closed, affecting function (chewing, speech) and/or aesthetics. Angle classification: Class I (neutrocclusion — normal molar relationship, but teeth may be crowded, rotated, spaced), Class II (retrognathism — mandible posterior to maxilla, overjet — typically >3-4 mm, Division 1: protruding maxillary incisors, Division 2: retroclined central incisors), Class III (prognathism — mandible anterior to maxilla, anterior crossbite, underbite). Other descriptors: overjet (horizontal overlap of maxillary incisors over mandibular, normal 2-3 mm), overbite (vertical overlap, normal 2-3 mm, deep bite >5 mm or covering >50% of lower incisors), open bite (anterior or posterior, no vertical overlap), crossbite (anterior or posterior, buccal or lingual), crowding, spacing, midline deviation, and impacted teeth. Etiology: genetic (skeletal discrepancies), environmental (thumb sucking, tongue thrust, mouth breathing, early loss of primary teeth), and acquired (periodontal disease, trauma). Treatment: orthodontic (braces, aligners, expanders, headgear), orthognathic surgery (for severe skeletal discrepancy), and restorative (crowns, veneers to correct minor malocclusion).

\medterm{Mesiodens}-- The most common supernumerary tooth (extra tooth), located in the midline of the maxilla between the central incisors, shaped like a peg or cone. Prevalence: 0.15-1.9%, more common in males (2:1). May be single or multiple, erupted or impacted. Complications: delayed eruption or impaction of permanent central incisors, midline diastema, root resorption of adjacent teeth, and cyst formation (dentigerous cyst). Most are asymptomatic and found incidentally on radiographs (occlusal or panoramic). Management: early removal (when roots of central incisors are 2/3 formed, around age 8-9) prevents orthodontic complications. Orthodontic closure of space after extraction. Associated with cleidocranial dysplasia and Gardner syndrome.

\medterm{Mucogingival Deformity}-- Condition involving the relationship between the gingiva and oral mucosa. Gingival recession (most common): apical migration of the gingival margin exposing the root surface, classified by Miller classification (Class I: not extending to mucogingival junction; Class II: extending to/beyond MGJ but no interproximal loss; Class III: interproximal loss; Class IV: severe interproximal loss). Causes: traumatic brushing, periodontal disease, tooth position (prominent root), thin biotype/fenestration, orthodontic movement, and frenum pull. Clinical: root exposure, dentin hypersensitivity (cold, air, touch), esthetic concern, and root caries risk. Treatment: non-surgical (desensitizing toothpaste, bonding/glass ionomer for root coverage), surgical root coverage (coronally advanced flap, subepithelial connective tissue graft — SCTG from the palate — gold standard for Class I and II, acellular dermal matrix allograft ADM, tunneling technique, and pinhole surgical technique). Free gingival graft (FGG): increases the zone of attached gingiva, not primarily for root coverage.

\medterm{Odontology}-- The scientific study of the structure, development, and diseases of the teeth and associated oral structures. Odontology encompasses dental anatomy, oral pathology, forensic odontology (identification of human remains through dental records, bite mark analysis), and clinical dentistry. Forensic odontology is particularly valuable when other identification methods fail, as teeth are the most durable structures in the human body and survive fire, decomposition, and trauma. Dental age estimation (through eruption patterns and cemental annulation) aids in both forensic identification and legal age determination.

\medterm{Oral and Maxillofacial Surgery (OMFS)}-- The dental specialty that diagnoses and surgically treats diseases, injuries, and defects of the mouth, jaws, face, and neck. Scope: dentoalveolar surgery (impacted third molar extraction — wisdom teeth, surgical extractions, pre-prosthetic surgery), dental implant placement, orthognathic surgery (corrective jaw surgery for malocclusion, skeletal discrepancies — Le Fort I maxillary osteotomy, bilateral sagittal split osteotomy BSSO for mandible), temporomandibular joint surgery (arthrocentesis, arthroscopy, discectomy, total joint replacement), facial trauma (mandible fractures, Le Fort II/III midface fractures, zygomaticomaxillary complex ZMC fractures, orbital floor blowout fractures, nasal fractures), oral pathology (biopsy, excision of lesions, management of oral cancer — composite resection, neck dissection, microvascular reconstruction), cleft lip and palate repair, and cosmetic facial surgery (rhinoplasty, facelift, blepharoplasty). Training: 4-6 year residency after dental school, often with an MD degree.

\medterm{Oral Medicine}-- The dental specialty focused on the oral health care of medically complex patients and the diagnosis and non-surgical management of oral mucosal and orofacial diseases. Areas: oral mucosal diseases (lichen planus, pemphigus vulgaris, mucous membrane pemphigoid, erythema multiforme, Sjogren syndrome — autoimmune exocrinopathy causing xerostomia and dry eyes), orofacial pain (temporomandibular disorders TMD, trigeminal neuralgia, burning mouth syndrome BMS, atypical odontalgia), salivary gland disorders (xerostomia, sialolithiasis, sialadenitis, Sjogren syndrome), oral complications of cancer therapy (mucositis, osteoradionecrosis, xerostomia, infection), and management of patients with systemic disease (diabetes, cardiovascular disease, bleeding disorders, bisphosphonate-related osteonecrosis of the jaw MRONJ).

\medterm{Oral Pathology}-- The dental specialty concerned with the diagnosis and management of diseases of the oral and maxillofacial region. Oral pathologists examine biopsies and provide histopathologic diagnosis. Common oral lesions: aphthous ulcers (canker sores — minor, major, herpetiform, recurrent aphthous stomatitis), herpes labialis (cold sores, HSV-1 reactivation), oral candidiasis (pseudomembranous thrush, erythematous, angular cheilitis), leukoplakia (white patch, cannot be scraped off, premalignant — 10-15% dysplastic), erythroplakia (red patch, higher malignant potential 50-70%), lichen planus (reticular white striae Wickham striae, erosive), oral lichenoid reaction, mucocele (mucus extravasation phenomenon), pyogenic granuloma (pregnancy tumor), fibroma, papilloma, and squamous cell carcinoma (most common oral malignancy, risk: tobacco, alcohol, HPV-16, betel nut).

\medterm{Oral Squamous Cell Carcinoma (OSCC)}-- The most common malignancy of the oral cavity, accounting for >90% of oral cancers. Risk factors: tobacco (smoking and smokeless — 75% of cases), alcohol (synergistic with tobacco), HPV-16 (oropharynx, base of tongue, tonsil — better prognosis), betel nut (Southeast Asia), immunosuppression (HIV, transplant), and sun exposure (lip). Premalignant lesions: leukoplakia (10-15% dysplastic or malignant), erythroplakia (50-70% high-risk), oral submucous fibrosis. Common sites: lateral tongue (most common), floor of mouth, buccal mucosa, gingiva, palate, and lip. Clinical: non-healing ulcer (>2 weeks), red/white patch, exophytic mass, pain, dysphagia, otalgia, loose teeth, cervical lymphadenopathy. Diagnosis: biopsy with histopathology, CT/MRI/PET-CT for staging. Treatment: surgical resection with negative margins (1 cm), neck dissection (therapeutic or elective for N0 with >15-20% occult metastasis risk), primary radiation (for early stage), adjuvant chemoradiation for high-risk features (positive margins, extracapsular nodal extension, perineural invasion). Prognosis: 60-70% 5-year survival for early stage, 30-50% for advanced.

\medterm{Orthodontics and Dentofacial Orthopedics}-- The dental specialty concerned with the diagnosis, prevention, and correction of malocclusion and dentofacial abnormalities. Orthodontists use fixed appliances (braces — metal, ceramic, lingual), clear aligners (Invisalign, Spark, ClearCorrect), and removable appliances to move teeth through bone by applying controlled forces. Key concepts: Angle classification of malocclusion — Class I (neutrocclusion, normal molar relationship, may have crowding/spacing), Class II (distocclusion, mandibular retrusion, overjet — Division 1 protruding maxillary incisors, Division 2 retroclined central incisors), Class III (mesiocclusion, mandibular prognathism, anterior crossbite). Cephalometric analysis: assesses skeletal and dental relationships on lateral cephalogram (SNA, SNB, ANB angles). Treatment phases: Phase I (interceptive, mixed dentition ages 7-10 — expander, partial braces, habit appliances), Phase II (comprehensive, full permanent dentition — braces/aligners, retention). Retention: fixed lingual retainers and/or clear vacuum-formed retainers to prevent relapse. Duration: 12-30 months. Complications: root resorption, white spot lesions (decalcification), gingival inflammation.

\medterm{Pediatric Dentistry (Pedodontics)}-- The dental specialty focusing on oral health care from infancy through adolescence, including children with special health care needs. Services: infant oral health exams (first visit by age 1), caries risk assessment, preventive care (fluoride varnish, sealants), restorative treatment for primary and young permanent teeth (composite fillings, stainless steel crowns for primary molars, pulp therapy), behavior management (tell-show-do, positive reinforcement, nitrous oxide sedation, protective stabilization, general anesthesia for extensive treatment), management of dental trauma (crown fractures, luxation, avulsion — reimplantation within 30 minutes best outcomes), space maintenance after early loss of primary teeth (band-and-loop, lingual arch, Nance appliance), and interceptive orthodontics (habit appliances — thumb/finger habit, tongue thrust). Early childhood caries (ECC): caries in primary teeth before age 6, severe if smooth-surface caries in children <3 years. Nursing/bottle caries: upper incisors affected by night-time bottle feeding.

\medterm{Periodontal (Gum) Disease}-- Also known as periodontal disease, an inflammatory condition affecting the gums and supporting structures of the teeth. It begins as gingivitis with redness, swelling, and bleeding of the gums and can progress to periodontitis, causing gum recession, bone loss, and tooth loss. Prevention and treatment involve good oral hygiene, professional dental cleanings, and in advanced cases, scaling, root planing, or surgery.

\medterm{Periodontal Abscess}-- A localized purulent infection in the periodontal pocket wall, usually an acute exacerbation of chronic periodontitis. Presents with ovoid, erythematous, fluctuant swelling on the gingiva along the lateral root surface, severe throbbing pain, tenderness to palpation/probing, tooth mobility, and purulent exudate on probing. The tooth is vital (unlike apical abscess where pulp is necrotic). Radiographic: advanced bone loss, often vertical or crater-shaped defects. Treatment: incision and drainage through the pocket, scaling and root planing (SRP), irrigation with chlorhexidine, occlusal adjustment if traumatic occlusion, and bite-block/tooth-stabilizing splint. Antibiotics if systemic involvement (amoxicillin with or without metronidazole, or clindamycin). Periodontal abscess must be distinguished from periapical abscess (pulp vitality testing, radiographic periapical vs. lateral lesion).

\medterm{Periodontal Disease}-- See Periodontics.

\medterm{Periodontics}-- The dental specialty concerned with the supporting structures of the teeth (periodontium: gingiva, periodontal ligament, cementum, alveolar bone), including diagnosis and treatment of periodontal disease. Gingivitis: reversible inflammation of the gingiva without bone loss, caused by bacterial plaque biofilm. Periodontitis: destructive inflammatory disease with progressive loss of alveolar bone and connective tissue attachment, classified by stage (I-IV, based on severity and complexity) and grade (A-C, based on rate of progression). Chronic periodontitis (most common): slow progression, adults >30. Aggressive periodontitis: rapid bone loss, younger patients, familial pattern. Necrotizing periodontal diseases: NUG (necrotizing ulcerative gingivitis — pain, fetor oris, pseudomembranes, fever) and NUP (necrotizing ulcerative periodontitis — associated with HIV, severe immunosuppression). Periodontal probing: 6-point pocket depth measurement (normal <3 mm), bleeding on probing (active inflammation), clinical attachment level (CAL). Non-surgical treatment: scaling and root planing (SRP, deep cleaning under local anesthesia). Surgical: flap surgery (open flap debridement), osseous surgery (recontour bone), guided tissue regeneration (GTR with barrier membrane), bone grafting, and platelet-rich plasma (PRP). Maintenance: periodontal maintenance q3-4 months in treated patients.

\medterm{Prosthodontics}-- The dental specialty concerned with the restoration and replacement of teeth using fixed and removable prostheses. Fixed prosthodontics: crowns (full coverage restoration for damaged teeth — metal, PFM, all-ceramic/e.max, zirconia), bridges (fixed partial denture replacing 1-2 missing teeth — traditional cantilever, resin-bonded Maryland bridge, implant-supported bridge), inlays and onlays (partial coverage restorations). Removable prosthodontics: complete dentures (full set for edentulous patients — conventional immediate, implant-retained overdenture), partial dentures (removable partial denture RPD with clasps, precision attachments). Implant prosthodontics: single crowns, implant-retained overdentures, implant-supported fixed complete dentures (All-on-4, All-on-6 concept using 4-6 implants with a full-arch fixed prosthesis). Maxillofacial prosthetics: obturator (for palatal defects after cancer surgery), speech aid prosthesis, nasal/auricular/ocular prosthesis, and radiation stent. Materials: polymethyl methacrylate (PMMA) denture base, cobalt-chromium alloy for RPD frameworks, ceramic (feldspathic, lithium disilicate, zirconia), and composite resin.

\medterm{Space Maintainer}-- A passive appliance used in pediatric dentistry to maintain the space created by the premature loss of a primary tooth, preventing adjacent teeth from drifting into the space and preserving sufficient space for the permanent successor to erupt. Types: band-and-loop (most common, for single tooth loss, unilateral — a band is cemented on the tooth adjacent to the space, with a loop extending across the edentulous space to contact the tooth on the other side), lingual arch (mandibular, maintains space when multiple primary molars are lost), Nance holding appliance (maxillary, similar to lingual arch with an acrylic button on the palate), and distal shoe (extends mesially into the extraction site to guide eruption of the first permanent molar). Crown and loop: for a missing primary molar when the adjacent tooth requires a full crown. Failure of space maintenance: space closure, crowding, impaction of the permanent tooth, and asymmetric eruption.

\medterm{Submandibular Space Infection}-- An infection of the submandibular space, usually from an odontogenic source (mandibular molar periapical infection spreading below the mylohyoid muscle). Presents with unilateral submandibular swelling, tenderness, and trismus. May progress to Ludwig angina (bilateral involvement with airway compromise). Imaging: CT with contrast (identifies involved spaces and abscess formation). Treatment: surgical drainage through the submandibular approach (incision parallel to the inferior border of the mandible, through platysma into the submandibular space, blunt dissection, Penrose drain), extraction of the causative tooth, and antibiotics (amoxicillin-clavulanate, clindamycin, or penicillin + metronidazole).

\medterm{Tartar (Dental Calculus)}-- Hardened, calcified dental plaque that forms on the teeth, both above the gumline (supragingival calculus) and below the gumline (subgingival calculus). Calculus is primarily composed of calcium phosphate (hydroxyapatite, brushite, whitlockite) and dead bacteria, mineralised from saliva (supragingival) or gingival crevicular fluid (subgingival). It is rough and porous, providing a surface for further plaque accumulation, and is a major risk factor for periodontal disease (gingivitis → periodontitis). Calculus cannot be removed by brushing; it requires professional scaling and root planing (dental cleaning) using ultrasonic scalers and hand instruments (curettes, scalers). Subgingival calculus appears brownish-black (stained by blood products). Tartar control toothpastes (pyrophosphates, zinc citrate) can reduce new calculus formation. Systemic risk factors: smoking (more calculus), xerostomia (less protective saliva), and diet (soft, processed foods).

\medterm{Teeth}-- See Tooth / Teeth and Dentition.

\medterm{Temporomandibular Joint Disorder (TMD)}-- A group of conditions causing pain and dysfunction in the temporomandibular joint, masticatory muscles, and associated structures. Etiology: multifactorial — malocclusion, bruxism (teeth grinding/clenching), trauma (whiplash, microtrauma), arthritis (osteoarthritis, rheumatoid, psoriatic joint disease), connective tissue disease, stress/anxiety (muscle hyperactivity), and parafunctional habits. Clinical: TMJ pain (preauricular, worse with jaw function), limited mouth opening (<35-40 mm), clicking/popping/crepitus (disc displacement with reduction = click on opening, without reduction = closed lock), locking (open or closed), muscle tenderness (masseter, temporalis, pterygoid), headaches, ear pain. Imaging: panoramic X-ray, CBCT for bone assessment, MRI for disc position and morphology. Treatment: conservative first-line (80% respond) — soft diet, NSAIDs, muscle relaxants, night guard/splint (flat-plane occlusal splint), physical therapy (jaw exercises, ultrasound, TENS), moist heat/ice. Second-line: arthrocentesis (joint lavage), arthroscopy (lysis and lavage, disc repositioning), and corticosteroid injection. Surgical: discectomy, total joint replacement (for severe degenerative disease). Botulinum toxin (Botox) for bruxism and masseter hypertrophy.

\medterm{Temporomandibular Joint Disorders (TMJ)}-- Temporomandibular joint disorders encompass a group of conditions affecting the jaw joint and the surrounding muscles that control jaw movement. Symptoms include jaw pain, clicking or popping sounds with movement, difficulty chewing, and locking of the jaw, often resulting from bruxism, arthritis, trauma, or malocclusion. Diagnosis is based on clinical examination and imaging when indicated; treatment ranges from conservative measures such as soft foods, stress reduction, and oral splints to physical therapy, medications, and in severe cases, arthrocentesis or surgery.

\medterm{TMJ Dislocation}-- Anterior displacement of the mandibular condyle from the glenoid fossa, usually occurring during wide mouth opening (yawning, dental treatment, vomiting, intubation). Presents with the mouth open, inability to close, preauricular depression (empty glenoid fossa), deviation of the mandible to the contralateral side, and pain. Diagnosis: clinical, confirmed by panoramic radiograph or CT. Reduction: the operator stands in front of the patient, places thumbs on the posterior mandibular molars (intraoral technique) or on the external oblique ridge, and applies downward and backward pressure to guide the condyle back into the fossa. Gauze between the thumbs to prevent being bitten. Muscle relaxants, sedation, or benzodiazepines may facilitate reduction. Post-reduction: soft diet, avoid yawning widely, NSAIDs. Recurrent dislocation: treated with autologous blood injection, capsular tightening, eminectomy, or zygomatic arch augmentation.

\medterm{Tooth Decay and Cavities}-- Tooth decay, also known as dental caries, is the progressive destruction of tooth enamel and dentin caused by acid produced by bacteria metabolizing sugars in the mouth. The process begins with demineralization of enamel and can progress through the dentin to the pulp, potentially causing pain, infection, and tooth loss. Prevention relies on regular brushing with fluoride toothpaste, flossing, limiting sugary foods and drinks, and professional applications of fluoride or dental sealants, while established cavities are treated with fillings, root canals, or extraction.

\medterm{Tooth Eruption}-- The process by which a tooth moves from its developmental position within the alveolar bone through the oral mucosa into functional occlusion. Primary (baby/deciduous) teeth: 20 teeth, eruption begins at 6 months (lower central incisors first), complete by 30 months. Order: central incisors (6-10 months), lateral incisors (8-12 months), first molars (12-16 months), canines (16-20 months), second molars (20-30 months). Permanent teeth: 32 teeth, begins at 6 years (first molar, lower central incisor), third molars (wisdom teeth ages 17-25). Timing variability is normal. Delayed eruption: caused by systemic conditions (hypothyroidism, hypoparathyroidism, rickets, cleidocranial dysplasia, Down syndrome) or local factors (over-retained primary tooth, supernumerary tooth, cyst, scar tissue from trauma). Eruption cyst: bluish soft tissue swelling over an erupting tooth in children; resolves spontaneously without treatment. Natal teeth: teeth present at birth (most common lower central incisors, may be mobile — extraction if risk of aspiration or injury during breastfeeding).

\medterm{Tooth Loss and Dental Implants}-- The absence of one or more teeth, most commonly resulting from advanced periodontal disease, severe tooth decay, or trauma, and can impair chewing, speech, and facial aesthetics. Dental implants are titanium posts surgically placed into the jawbone that serve as artificial tooth roots, providing a stable foundation for crowns, bridges, or dentures. Implant-supported restorations offer superior function and longevity compared to traditional removable prosthetics, though they require adequate bone density, good oral hygiene, and a healing period of several months.

\medterm{Toothache (Odontalgia)}-- Pain originating from the tooth or its supporting structures. This is the most common reason for emergency dental visits. Common causes: dental caries (most common — deep caries extending to dentin/pulp), pulpitis (reversible: sharp, brief pain to cold/sweet, resolves quickly; irreversible: spontaneous, lingering, throbbing pain to cold, may last minutes), abscess (periapical or periodontal — severe pain, tenderness to biting/percussion), cracked tooth syndrome (sharp, fleeting pain on biting/release — may be intermittent, difficult to diagnose), dry socket (alveolar osteitis — severe post-extraction pain with referral to ear, 2-3 days post-op), pericoronitis (pain/swelling around erupting third molar), sinusitis (maxillary — referred pain to upper teeth, pain on bending forward), and trigeminal neuralgia (electric-shock unilateral pain triggered by light touch). Diagnosis: history (character of pain, duration, triggers), clinical exam (caries probing, percussion, palpation, mobility, periodontal probing), vitality testing (cold test — refrigerant spray, electric pulp tester — or heat test for diagnosis of irreversible pulpitis), and radiographs (bitewing, periapical, panoramic, CBCT). Treatment: depends on diagnosis — restoration (filling/crown), pulp capping, root canal therapy, extraction, periodontal treatment, or anti-inflammatory/muscle relaxant for TMD.

\medterm{Wisdom Tooth (Third Molar)}-- The third molar, the last tooth in the dental arch, typically erupting between ages 17 and 25 (hence wisdom tooth). Humans have four third molars (maxillary right and left, mandibular right and left). Wisdom teeth are the most commonly impacted (unerupted or partially erupted) teeth due to insufficient space in the jaw, with an impaction rate of 25-75%. Impaction types: mesioangular (most common, angled forward), distoangular (angled backward), vertical (upright but blocked), horizontal (lying on side, highest risk of mandibular nerve injury during extraction), and buccal/lingual. Complications: pericoronitis (inflammation and infection of the soft tissue around a partially erupted wisdom tooth, presenting with pain, trismus, swelling, halitosis, and purulent discharge), caries (distal surface of second molar), periodontal disease, cyst formation (dentigerous cyst around the impacted crown), root resorption of adjacent teeth, and rarely tumour (ameloblastoma, odontogenic keratocyst). Management: prophylactic extraction of asymptomatic, disease-free third molars is controversial; NICE guidelines recommend against prophylactic extraction of healthy, non-pathological wisdom teeth. Indications for extraction: recurrent pericoronitis (2+ episodes), caries or periodontal disease not amenable to restoration, pathology (cyst, tumour), external root resorption of the second molar, and for orthodontic reasons (lack of space, relapse prevention). Surgical extraction: performed under local anaesthesia, intravenous sedation, or general anaesthesia. Assessment: panoramic radiograph (orthopantomogram, OPG) to evaluate root morphology, number of roots (1-4), root curvature, proximity to the inferior alveolar nerve canal (mandibular canal). Complications of extraction: dry socket (alveolar osteitis, 5-10% of mandibular extractions), pain, swelling, trismus, bleeding, infection, and nerve injury (inferior alveolar nerve — lip paraesthesia/anaesthesia, incidence 0.5-5%; lingual nerve — tongue paraesthesia, incidence 0.1-2%). Risk of nerve injury is increased with horizontal and distoangular impactions, deep impaction, and when the nerve canal is superimposed on the root tips on OPG.
